\subtituloA{Multidisciplinariedad del desarrollo web}

\lista{
  \textbf{Tecnologías empleadas:} La elección de lenguajes, frameworks y herramientas define las posibilidades y limitaciones del proyecto. |
  \textbf{Distribución y rendimiento:} La arquitectura del sitio y su capacidad de respuesta impactan directamente en la usabilidad. |
  \textbf{Propósito del sitio:} Cada web tiene un objetivo específico, ya sea comercial, informativo, educativo o de otro tipo.
}

\nota{Para el examen}{
¿Cuál es el objetivo principal del World Wide Web Consortium (W3C)?
a) Promover la creación de navegadores alternativos
b) Establecer estándares técnicos para garantizar la interoperabilidad web
c) Proveer servicios de alojamiento web
d) Desarrollar lenguajes de programación para servidores. La

\textbf{Respuesta:} es b), ya que el W3C crea y mantiene especificaciones técnicas que permiten que diferentes navegadores interpreten de manera coherente los contenidos web. La opción
a) es incorrecta porque el W3C no desarrolla navegadores. La opción
c) es falsa porque el alojamiento web es responsabilidad de los proveedores de servicios. La opción
d) es incorrecta porque los lenguajes de programación son creados por sus comunidades, no por el W3C.
}